\documentclass[aspectratio=169]{beamer}
% Add slides directory to search path for Dartmouth theme
\makeatletter
\def\input@path{{../}}
\makeatother
\usetheme{Dartmouth}
\usepackage{graphicx}
\usepackage{tikz}
\usepackage{amsmath}
\usepackage{listings}
\usepackage{xcolor}
\usepackage{booktabs}
\usetikzlibrary{shapes,arrows,positioning,calc}

% Code listing setup
\lstset{
  basicstyle=\ttfamily\small,
  keywordstyle=\color{blue},
  commentstyle=\color{gray},
  stringstyle=\color{red},
  showstringspaces=false,
  breaklines=true,
  frame=single,
  language=Python
}

\title{Lecture 17: Applications of Encoder Models}
\subtitle{Week 6, Lecture 3 - From Theory to Practice}
\author{PSYC 51.07: Models of Language and Communication}
\date{Winter 2026}

\begin{document}

% ============================================================
% TITLE SLIDE
% ============================================================
\begin{frame}[shrink=10]
\titlepage
\end{frame}

% ============================================================
% OVERVIEW
% ============================================================
\begin{frame}[shrink=10]{Today's Agenda 📋}
\large
\begin{enumerate}\setlength{\itemsep}{0pt}
    \item 🚀 \textbf{Real-World Applications}: Where BERT shines
    \item 🧠 \textbf{Cognitive Neuroscience}: Brain-model parallels
    \item 🤔 \textbf{Understanding vs. Pattern Matching}: The big debate
    \item ⚠️ \textbf{Limitations}: What BERT can't do
    \item 💡 \textbf{Practical Tips}: Deployment and optimization
    \item 🔮 \textbf{Future Directions}: Where are we heading?
\end{enumerate}

\vspace{1em}
\textit{Goal: Connect BERT to real applications and understand broader implications}
\end{frame}

% ============================================================
% BERT APPLICATIONS
% ============================================================
\begin{frame}[shrink=10]{BERT Applications 🚀}
\small
\textbf{BERT excels at understanding tasks:}

\vspace{1em}

\begin{columns}[T]
\begin{column}{0.48\textwidth}
\textbf{Classification Tasks:}
\begin{itemize}\setlength{\itemsep}{0pt}
    \item Sentiment Analysis
    \item Topic Classification
    \item Spam Detection
    \item Intent Recognition
\end{itemize}

\vspace{0.5em}

\textbf{Token-Level Tasks:}
\begin{itemize}\setlength{\itemsep}{0pt}
    \item Named Entity Recognition (NER)
    \item Part-of-Speech Tagging
    \item Word Sense Disambiguation
\end{itemize}
\end{column}

\begin{column}{0.48\textwidth}
\textbf{Span-Level Tasks:}
\begin{itemize}\setlength{\itemsep}{0pt}
    \item Question Answering
    \item Extractive Summarization
    \item Information Extraction
\end{itemize}

\vspace{0.5em}

\textbf{Sentence-Pair Tasks:}
\begin{itemize}\setlength{\itemsep}{0pt}
    \item Semantic Similarity
    \item Natural Language Inference
    \item Paraphrase Detection
\end{itemize}
\end{column}
\end{columns}

\vspace{1em}

\begin{block}{Industry Impact}
BERT powers:
\begin{itemize}\setlength{\itemsep}{0pt}
    \item Google Search (understanding queries)
    \item Customer service chatbots
    \item Content moderation
    \item Document understanding
\end{itemize}
\end{block}
\end{frame}

% ============================================================
% GOOGLE SEARCH
% ============================================================
\begin{frame}[shrink=10]{Case Study: Google Search 🔍}
\textbf{BERT revolutionized search in 2019}

\vspace{1em}

\textbf{Problem:}
\begin{itemize}\setlength{\itemsep}{0pt}
    \item Search queries often depend on context and word order
    \item Traditional keyword matching misses nuance
    \item Example: "2019 brazil traveler to usa need a visa"
\end{itemize}

\vspace{1em}

\begin{columns}[T]
\begin{column}{0.48\textwidth}
\textbf{Before BERT:}
\begin{itemize}\setlength{\itemsep}{0pt}
    \item Focused on "brazil" and "usa"
    \item Missed importance of "to"
    \item Returned results about US travelers to Brazil
    \item \alert{Wrong direction!}
\end{itemize}

\vspace{0.5em}

\begin{center}
\includegraphics[width=0.9\textwidth]{example-image-a}

\small Wrong: US → Brazil
\end{center}
\end{column}

\begin{column}{0.48\textwidth}
\textbf{With BERT:}
\begin{itemize}\setlength{\itemsep}{0pt}
    \item Understands "to" is critical
    \item Grasps full context and direction
    \item Returns correct results about Brazilian travelers to US
    \item \alert{Correct understanding!}
\end{itemize}

\vspace{0.5em}

\begin{center}
\includegraphics[width=0.9\textwidth]{example-image-b}

\small Correct: Brazil → US
\end{center}
\end{column}
\end{columns}

\vspace{0.5em}
\small
\textit{Google reported BERT improved 1 in 10 searches in English}
\end{frame}

% ============================================================
% QA EXAMPLE
% ============================================================
\begin{frame}[fragile]{Question Answering with BERT 💬}
\textbf{Extractive QA: Find answer span in passage}

\vspace{1em}

\begin{exampleblock}{Example}
\textbf{Context:} "The Normans (Norman: Nourmands; French: Normands; Latin: Normanni) were the people who in the 10th and 11th centuries gave their name to Normandy, a region in France."

\vspace{0.3em}

\textbf{Question:} "In what country is Normandy located?"

\vspace{0.3em}

\textbf{Answer:} France
\end{exampleblock}

\vspace{1em}

\begin{lstlisting}[language=Python]
from transformers import pipeline

# Load QA pipeline with BERT
qa_pipeline = pipeline("question-answering", model="bert-large-uncased-whole-word-masking-finetuned-squad")

# Ask question
result = qa_pipeline(
    question="In what country is Normandy located?",
    context="The Normans were the people who in the 10th and 11th centuries gave their name to Normandy, a region in France."
)

print(result)
# {'answer': 'France', 'score': 0.987, 'start': 134, 'end': 140}
\end{lstlisting}

\vspace{0.5em}
\textbf{BERT predicts start and end positions of the answer span!}
\end{frame}

% ============================================================
% NER EXAMPLE
% ============================================================
\begin{frame}[fragile]{Named Entity Recognition 🏷️}
\textbf{Token-level classification task}

\vspace{1em}

\begin{exampleblock}{Example}
\textbf{Input:} "Apple Inc. is headquartered in Cupertino, California."

\vspace{0.3em}

\textbf{Output:}
\begin{itemize}\setlength{\itemsep}{0pt}
    \item Apple Inc. → \textcolor{blue}{ORGANIZATION}
    \item Cupertino → \textcolor{green}{LOCATION}
    \item California → \textcolor{green}{LOCATION}
\end{itemize}
\end{exampleblock}

\vspace{1em}

\begin{lstlisting}[language=Python]
from transformers import pipeline

# Load NER pipeline
ner_pipeline = pipeline("ner", model="dslim/bert-base-NER")

# Extract entities
text = "Apple Inc. is headquartered in Cupertino, California."
entities = ner_pipeline(text)

for entity in entities:
    print(f"{entity['word']}: {entity['entity']} (score: {entity['score']:.2f})")

# Output:
# Apple: B-ORG (score: 0.99)
# Inc: I-ORG (score: 0.99)
# Cupertino: B-LOC (score: 0.99)
# California: B-LOC (score: 0.99)
\end{lstlisting}
\end{frame}

% ============================================================
% SENTIMENT ANALYSIS
% ============================================================
\begin{frame}[fragile]{Sentiment Analysis 😊😐😢}
\textbf{Sequence classification task}

\vspace{1em}

\begin{exampleblock}{Examples}
\begin{itemize}\setlength{\itemsep}{0pt}
    \item "This movie was absolutely amazing!" → \textcolor{green}{POSITIVE}
    \item "The product broke after one week." → \textcolor{red}{NEGATIVE}
    \item "The weather is cloudy today." → \textcolor{gray}{NEUTRAL}
\end{itemize}
\end{exampleblock}

\vspace{1em}

\begin{lstlisting}[language=Python]
from transformers import pipeline

# Load sentiment analysis pipeline
sentiment_pipeline = pipeline("sentiment-analysis", model="distilbert-base-uncased-finetuned-sst-2-english")

# Analyze sentiments
texts = [
    "This movie was absolutely amazing!",
    "The product broke after one week.",
    "The weather is cloudy today."
]

for text in texts:
    result = sentiment_pipeline(text)[0]
    print(f"{text}")
    print(f"  → {result['label']} (confidence: {result['score']:.2f})\n")
\end{lstlisting}

\vspace{0.5em}
\textbf{Applications:} Customer reviews, social media monitoring, brand sentiment
\end{frame}

% ============================================================
% SEMANTIC SIMILARITY
% ============================================================
\begin{frame}[fragile]{Semantic Similarity 🔗}
\textbf{Measuring sentence similarity with BERT embeddings}

\vspace{1em}

\begin{lstlisting}[language=Python]
from transformers import BertTokenizer, BertModel
import torch
import torch.nn.functional as F

tokenizer = BertTokenizer.from_pretrained('bert-base-uncased')
model = BertModel.from_pretrained('bert-base-uncased')

def get_sentence_embedding(sentence):
    inputs = tokenizer(sentence, return_tensors='pt', padding=True, truncation=True)
    outputs = model(**inputs)
    # Use [CLS] token embedding as sentence representation
    return outputs.last_hidden_state[:, 0, :]

# Compare sentences
sent1 = "The cat is sleeping on the couch"
sent2 = "A feline is resting on the sofa"
sent3 = "The weather is nice today"

emb1 = get_sentence_embedding(sent1)
emb2 = get_sentence_embedding(sent2)
emb3 = get_sentence_embedding(sent3)

# Compute cosine similarities
sim_12 = F.cosine_similarity(emb1, emb2).item()
sim_13 = F.cosine_similarity(emb1, emb3).item()

print(f"Similarity (1-2): {sim_12:.3f}")  # High (paraphrases)
print(f"Similarity (1-3): {sim_13:.3f}")  # Low (different topics)
\end{lstlisting}
\end{frame}

% ============================================================
% COGNITIVE NEUROSCIENCE
% ============================================================
\begin{frame}[shrink=10]{Cognitive Neuroscience Perspective 🧠}
\textbf{How do brains and models process language?}

\vspace{1em}

\begin{columns}[T]
\begin{column}{0.48\textwidth}
\textbf{Predictive Processing in the Brain:}
\begin{itemize}\setlength{\itemsep}{0pt}
    \item Brain constantly predicts upcoming input
    \item N400: Neural response to unexpected words
    \item P600: Syntactic anomaly detection
    \item Context shapes predictions
    \item Prediction errors drive learning
\end{itemize}

\vspace{0.5em}

\textbf{Key Brain Regions:}
\begin{itemize}\setlength{\itemsep}{0pt}
    \item \textbf{Left IFG}: Syntax processing
    \item \textbf{Left STG/MTG}: Semantic processing
    \item \textbf{ATL}: Conceptual knowledge
\end{itemize}
\end{column}

\begin{column}{0.48\textwidth}
\textbf{Predictive Processing in Models:}
\begin{itemize}\setlength{\itemsep}{0pt}
    \item BERT: Predict masked words
    \item GPT: Predict next word
    \item Both use context to predict
    \item Surprise = high loss
    \item Gradient descent = learning
\end{itemize}

\vspace{0.5em}

\textbf{Similarities:}
\begin{itemize}\setlength{\itemsep}{0pt}
    \item Both hierarchical
    \item Both context-sensitive
    \item Both predictive
    \item Both learn from errors
\end{itemize}
\end{column}
\end{columns}

\vspace{0.5em}
\small
\textit{References: Kuperberg \& Jaeger (2016), Willems et al. (2016), Hagoort \& Indefrey (2014)}
\end{frame}

% ============================================================
% PREDICTION IN BRAINS VS MODELS
% ============================================================
\begin{frame}{Prediction in Brains vs. Language Models 🧠🤖}
\textbf{Parallels between neural and artificial systems}

\vspace{1em}

\begin{center}
\begin{tabular}{p{5cm}p{5cm}}
\toprule
\textbf{Human Brain} & \textbf{Transformer Models} \\
\midrule
N400 component\\
(surprise at unexpected word) &
High cross-entropy loss\\
(low probability prediction) \\
\midrule
Hierarchical processing\\
(sounds → words → sentences) &
Layer-by-layer representations\\
(tokens → phrases → meaning) \\
\midrule
Context integration\\
(prior discourse, world knowledge) &
Self-attention mechanism\\
(attend to relevant context) \\
\midrule
Prediction error signals &
Gradient-based learning \\
\midrule
Distributed representations\\
(population coding) &
Distributed embeddings\\
(vector representations) \\
\bottomrule
\end{tabular}
\end{center}

\vspace{1em}

\textbf{Question:} Are these superficial analogies or deep connections?

\small
\textit{Reference: Kuperberg \& Jaeger (2016) - "What do we mean by prediction in language comprehension?"}
\end{frame}

% ============================================================
% NEURAL ENCODING
% ============================================================
\begin{frame}[shrink=10]{Neural Encoding with Language Models 🔬}
\textbf{Can we predict brain activity from language models?}

\vspace{1em}

\textbf{Approach:}
\begin{enumerate}\setlength{\itemsep}{0pt}
    \item Show participants text while recording brain activity (fMRI/EEG)
    \item Extract representations from language model (e.g., BERT layer 6)
    \item Train regression model: Model embeddings → Brain activity
    \item Test: Can we predict brain responses to new text?
\end{enumerate}

\vspace{1em}

\textbf{Findings:}
\begin{itemize}\setlength{\itemsep}{0pt}
    \item \textbf{Better models → Better brain prediction}
    \item BERT outperforms Word2Vec at predicting brain activity
    \item Different layers correlate with different brain regions
    \item Middle layers best for semantic areas
    \item Suggests shared representations?
\end{itemize}

\vspace{1em}

\begin{alertblock}{But...}
Correlation ≠ Causation! Models might capture statistical regularities that brains use, without using the same mechanisms.
\end{alertblock}

\small
\textit{Reference: Willems et al. (2016) - "Prediction during natural language comprehension"}
\end{frame}

% ============================================================
% WHAT DOES MODEL UNDERSTAND
% ============================================================
\begin{frame}[shrink=10]{Discussion: What Does the Model "Understand"? 🤔}
\small
\begin{center}
\LARGE
\textbf{Does BERT understand language?}
\end{center}

\vspace{1em}

\begin{columns}[T]
\begin{column}{0.48\textwidth}
\textbf{Evidence FOR understanding:}
\begin{itemize}\setlength{\itemsep}{0pt}
    \item Captures syntax and semantics
    \item Resolves ambiguity
    \item Handles long-range dependencies
    \item Generalizes to new examples
    \item Predicts brain activity
    \item Solves complex tasks
\end{itemize}

\vspace{0.5em}
\textit{"If it acts like it understands, maybe it does?"}
\end{column}

\begin{column}{0.48\textwidth}
\textbf{Evidence AGAINST understanding:}
\begin{itemize}\setlength{\itemsep}{0pt}
    \item No grounding in physical world
    \item No sensory experience
    \item No social context
    \item Brittle to adversarial examples
    \item No common sense reasoning
    \item Only pattern matching?
\end{itemize}

\vspace{0.5em}
\textit{"Understanding requires more than statistical patterns"}
\end{column}
\end{columns}

\vspace{1em}

\begin{block}{Key Questions}
\begin{itemize}\setlength{\itemsep}{0pt}
    \item What is the difference between understanding and correlation?
    \item Can meaning exist without grounding?
    \item Is human understanding fundamentally different?
\end{itemize}
\end{block}

\textbf{Class Discussion: What do YOU think?}
\end{frame}

% ============================================================
% ADVERSARIAL EXAMPLES
% ============================================================
\begin{frame}{Adversarial Examples and Brittleness ⚠️}
\textbf{BERT can be fooled easily}

\vspace{1em}

\begin{exampleblock}{Example: Sentiment Analysis}
\textbf{Original:} "This movie was great!"

\textbf{Prediction:} POSITIVE ✓

\vspace{0.5em}

\textbf{Modified:} "This movie was great! [SEP] terrible terrible terrible terrible"

\textbf{Prediction:} NEGATIVE ✗

\vspace{0.5em}

Just adding repeated negative words after [SEP] flips the prediction!
\end{exampleblock}

\vspace{1em}

\begin{alertblock}{Other Brittleness Issues}
\begin{itemize}\setlength{\itemsep}{0pt}
    \item \textbf{Word substitutions}: Replace "good" with synonym "excellent" → wrong prediction
    \item \textbf{Paraphrases}: Semantically equivalent but different words → different predictions
    \item \textbf{Meaningless additions}: Add irrelevant text → changes prediction
    \item \textbf{Typos and noise}: Small perturbations cause large errors
\end{itemize}
\end{alertblock}

\vspace{0.5em}

\textbf{Implication:} Models learn spurious patterns, not robust understanding
\end{frame}

% ============================================================
% LIMITATIONS
% ============================================================
\begin{frame}[shrink=10]{Limitations of Current Models ⚠️}
\textbf{Despite impressive performance, transformers have limitations:}

\vspace{1em}

\begin{enumerate}\setlength{\itemsep}{0pt}
    \item \textbf{Quadratic Complexity}
    \begin{itemize}\setlength{\itemsep}{0pt}
        \item Self-attention scales as $O(n^2)$
        \item Limited context windows (512-4096 tokens)
        \item Cannot process very long documents efficiently
    \end{itemize}

    \item \textbf{No True Understanding}
    \begin{itemize}\setlength{\itemsep}{0pt}
        \item Pattern matching vs. comprehension
        \item Lack of common sense
        \item No world model
    \end{itemize}

    \item \textbf{Data Efficiency}
    \begin{itemize}\setlength{\itemsep}{0pt}
        \item Requires massive training data
        \item Humans learn language with much less data
        \item Not biologically plausible
    \end{itemize}

    \item \textbf{Biases and Fairness}
    \begin{itemize}\setlength{\itemsep}{0pt}
        \item Inherits biases from training data
        \item Can amplify stereotypes
        \item Ethical concerns
    \end{itemize}
\end{enumerate}
\end{frame}

% ============================================================
% BIAS EXAMPLE
% ============================================================
\begin{frame}[shrink=10]{Bias in Language Models ⚖️}
\textbf{Models reflect and can amplify societal biases}

\vspace{1em}

\begin{alertblock}{Examples of Bias}
\begin{itemize}\setlength{\itemsep}{0pt}
    \item \textbf{Gender bias:}
    \begin{itemize}\setlength{\itemsep}{0pt}
        \item "The doctor said [MASK] was running late" → he (90\%)
        \item "The nurse said [MASK] was running late" → she (80\%)
    \end{itemize}

    \item \textbf{Occupational stereotypes:}
    \begin{itemize}\setlength{\itemsep}{0pt}
        \item Associates certain jobs with certain genders
        \item "Programmer" → male pronouns
        \item "Secretary" → female pronouns
    \end{itemize}

    \item \textbf{Racial bias:}
    \begin{itemize}\setlength{\itemsep}{0pt}
        \item Different sentiment for names associated with different races
        \item Can perpetuate harmful stereotypes
    \end{itemize}
\end{itemize}
\end{alertblock}

\vspace{1em}

\textbf{Mitigation Strategies:}
\begin{itemize}\setlength{\itemsep}{0pt}
    \item Debiasing techniques during training
    \item Balanced and diverse training data
    \item Post-processing and filtering
    \item Human oversight and auditing
    \item Ongoing research area!
\end{itemize}
\end{frame}

% ============================================================
% PRACTICAL TIPS
% ============================================================
\begin{frame}[shrink=10]{Practical Tips for Working with Transformers 💡}
\begin{enumerate}\setlength{\itemsep}{0pt}
    \item \textbf{Start with Pre-trained Models}
    \begin{itemize}\setlength{\itemsep}{0pt}
        \item Don't train from scratch (too expensive!)
        \item Use HuggingFace Model Hub
        \item Choose appropriate model size
    \end{itemize}

    \item \textbf{Fine-tuning Best Practices}
    \begin{itemize}\setlength{\itemsep}{0pt}
        \item Use small learning rate (1e-5 to 5e-5)
        \item Add warmup steps
        \item Monitor for overfitting
        \item Freeze early layers if data is limited
    \end{itemize}

    \item \textbf{Computational Efficiency}
    \begin{itemize}\setlength{\itemsep}{0pt}
        \item Use mixed precision training (FP16)
        \item Gradient accumulation for larger batch sizes
        \item Consider DistilBERT for faster inference
        \item Use FlashAttention when available
    \end{itemize}

    \item \textbf{Evaluation}
    \begin{itemize}\setlength{\itemsep}{0pt}
        \item Use task-specific metrics
        \item Test on out-of-distribution data
        \item Check for biases
        \item Visualize attention for interpretability
    \end{itemize}
\end{enumerate}
\end{frame}

% ============================================================
% DEPLOYMENT
% ============================================================
\begin{frame}[shrink=10]{Deployment Considerations 🚀}
\textbf{Moving from research to production}

\vspace{1em}

\begin{enumerate}\setlength{\itemsep}{0pt}
    \item \textbf{Model Selection}
    \begin{itemize}\setlength{\itemsep}{0pt}
        \item Balance quality vs. speed
        \item DistilBERT for low-latency applications
        \item BERT-Base for balanced performance
        \item BERT-Large when quality is critical
    \end{itemize}

    \vspace{0.5em}

    \item \textbf{Optimization}
    \begin{itemize}\setlength{\itemsep}{0pt}
        \item Quantization: FP32 → INT8 (4x smaller, faster)
        \item Model pruning: Remove unnecessary weights
        \item Knowledge distillation: Compress to smaller model
        \item ONNX Runtime for optimized inference
    \end{itemize}

    \vspace{0.5em}

    \item \textbf{Infrastructure}
    \begin{itemize}\setlength{\itemsep}{0pt}
        \item GPU acceleration for batch processing
        \item CPU optimization for single requests
        \item Caching for repeated queries
        \item Load balancing for high traffic
    \end{itemize}

    \vspace{0.5em}

    \item \textbf{Monitoring}
    \begin{itemize}\setlength{\itemsep}{0pt}
        \item Track latency and throughput
        \item Monitor prediction quality
        \item Detect distribution drift
        \item A/B testing for improvements
    \end{itemize}
\end{enumerate}
\end{frame}

% ============================================================
% FUTURE DIRECTIONS
% ============================================================
\begin{frame}[shrink=10]{Future Directions 🔮}
\textbf{Where is the field heading?}

\vspace{1em}

\begin{enumerate}\setlength{\itemsep}{0pt}
    \item \textbf{Longer Context}
    \begin{itemize}\setlength{\itemsep}{0pt}
        \item Efficient attention mechanisms (linear, sparse)
        \item Models with 100K+ token context
        \item Better long-document understanding
    \end{itemize}

    \vspace{0.5em}

    \item \textbf{Multimodal Models}
    \begin{itemize}\setlength{\itemsep}{0pt}
        \item Vision + Language (CLIP, DALL-E)
        \item Audio + Language (Whisper)
        \item Grounded understanding
    \end{itemize}

    \vspace{0.5em}

    \item \textbf{Better Pre-training}
    \begin{itemize}\setlength{\itemsep}{0pt}
        \item More efficient objectives
        \item Curriculum learning
        \item Continual learning
    \end{itemize}

    \vspace{0.5em}

    \item \textbf{Smaller, More Efficient Models}
    \begin{itemize}\setlength{\itemsep}{0pt}
        \item Better compression techniques
        \item Lottery ticket hypothesis
        \item Edge deployment
    \end{itemize}

    \vspace{0.5em}

    \item \textbf{Addressing Limitations}
    \begin{itemize}\setlength{\itemsep}{0pt}
        \item Debiasing and fairness
        \item Robustness and adversarial training
        \item Common sense reasoning
        \item Interpretability and explainability
    \end{itemize}
\end{enumerate}
\end{frame}

% ============================================================
% ENCODER VS DECODER REVISITED
% ============================================================
\begin{frame}{Encoder vs Decoder Models Revisited 🔄}
\textbf{Different models for different tasks}

\vspace{1em}

\begin{center}
\begin{tabular}{lp{4.5cm}p{4.5cm}}
\toprule
 & \textbf{Encoder (BERT)} & \textbf{Decoder (GPT)} \\
\midrule
\textbf{Architecture} &
Bidirectional self-attention &
Unidirectional (causal) attention \\
\midrule
\textbf{Pre-training} &
Masked Language Modeling &
Next token prediction \\
\midrule
\textbf{Best for} &
Understanding tasks:\\
• Classification\\
• NER, QA\\
• Similarity &
Generation tasks:\\
• Text completion\\
• Dialogue\\
• Creative writing \\
\midrule
\textbf{Trend} &
Dominant 2018-2020 &
Increasingly popular 2020+ \\
\bottomrule
\end{tabular}
\end{center}

\vspace{1em}

\textbf{Interesting observation:}
\begin{itemize}\setlength{\itemsep}{0pt}
    \item Decoder-only models (GPT-3, LLaMA) can also do classification via prompting!
    \item "In-context learning" blurs the distinction
    \item Trend toward unified decoder-only architectures
\end{itemize}
\end{frame}

% ============================================================
% DISCUSSION
% ============================================================
\begin{frame}[shrink=10]{Discussion Questions 💭}
\begin{enumerate}\setlength{\itemsep}{0pt}
    \item \textbf{Understanding vs. Pattern Matching:}
    \begin{itemize}\setlength{\itemsep}{0pt}
        \item Where do you draw the line?
        \item Is there a test for "true" understanding?
        \item Does it matter for applications?
    \end{itemize}

    \vspace{0.5em}

    \item \textbf{Brain-Model Parallels:}
    \begin{itemize}\setlength{\itemsep}{0pt}
        \item How useful are these comparisons?
        \item What can neuroscience learn from AI?
        \item What can AI learn from neuroscience?
    \end{itemize}

    \vspace{0.5em}

    \item \textbf{Bias and Fairness:}
    \begin{itemize}\setlength{\itemsep}{0pt}
        \item Who is responsible for addressing bias?
        \item Can we ever have completely unbiased models?
        \item How do we balance accuracy and fairness?
    \end{itemize}

    \vspace{0.5em}

    \item \textbf{Future of NLP:}
    \begin{itemize}\setlength{\itemsep}{0pt}
        \item Will encoder models remain relevant?
        \item Are decoder-only models the future?
        \item What's the next big breakthrough?
    \end{itemize}
\end{enumerate}
\end{frame}

% ============================================================
% ASSIGNMENT 4
% ============================================================
\begin{frame}[shrink=10]{Assignment 4: Context-Aware Models 📝}
\textbf{Hands-on experience with transformers!}

\vspace{1em}

\textbf{Tasks:}
\begin{enumerate}\setlength{\itemsep}{0pt}
    \item \textbf{Implement Attention Mechanism}
    \begin{itemize}\setlength{\itemsep}{0pt}
        \item Build scaled dot-product attention from scratch
        \item Visualize attention weights
    \end{itemize}

    \item \textbf{Fine-tune BERT}
    \begin{itemize}\setlength{\itemsep}{0pt}
        \item Load pre-trained BERT
        \item Fine-tune on sentiment analysis
        \item Compare to baseline models
    \end{itemize}

    \item \textbf{Analyze Contextual Embeddings}
    \begin{itemize}\setlength{\itemsep}{0pt}
        \item Extract embeddings for polysemous words
        \item Visualize how context changes representations
        \item Compare BERT vs Word2Vec
    \end{itemize}

    \item \textbf{Explore Different Architectures}
    \begin{itemize}\setlength{\itemsep}{0pt}
        \item Compare BERT (encoder) vs GPT (decoder)
        \item Test on different tasks
        \item Analyze strengths/weaknesses
    \end{itemize}

    \item \textbf{Research Component}
    \begin{itemize}\setlength{\itemsep}{0pt}
        \item Read one paper from references
        \item Write brief summary \& critical analysis
    \end{itemize}
\end{enumerate}

\vspace{0.5em}
\textbf{Due:} Check course website for deadline
\end{frame}

% ============================================================
% SUMMARY WEEKS 5-6
% ============================================================
\begin{frame}[shrink=10]{Summary: Weeks 5-6 🎯}
\textbf{What we learned:}

\vspace{1em}

\begin{enumerate}\setlength{\itemsep}{0pt}
    \item \textbf{Evolution of Context}
    \begin{itemize}\setlength{\itemsep}{0pt}
        \item Seq2Seq → Attention → Transformers
        \item From bottleneck to full parallelization
    \end{itemize}

    \item \textbf{Transformer Architecture}
    \begin{itemize}\setlength{\itemsep}{0pt}
        \item Self-attention, multi-head attention
        \item Positional encoding, layer norm, residuals
        \item Encoder-only (BERT), Decoder-only (GPT), Both (T5)
    \end{itemize}

    \item \textbf{BERT \& Variants}
    \begin{itemize}\setlength{\itemsep}{0pt}
        \item Masked Language Modeling
        \item Pre-train then fine-tune paradigm
        \item RoBERTa, ALBERT, DistilBERT, ELECTRA
    \end{itemize}

    \item \textbf{Applications}
    \begin{itemize}\setlength{\itemsep}{0pt}
        \item Classification, NER, QA, similarity
        \item Real-world impact (Google Search, etc.)
    \end{itemize}

    \item \textbf{Broader Implications}
    \begin{itemize}\setlength{\itemsep}{0pt}
        \item Brain-model parallels
        \item Understanding vs. pattern matching
        \item Limitations and future directions
    \end{itemize}
\end{enumerate}
\end{frame}

% ============================================================
% RESOURCES
% ============================================================
\begin{frame}[shrink=10]{Resources \& Further Reading 📚}
\textbf{Key Papers:}
\begin{itemize}\setlength{\itemsep}{0pt}
    \item Vaswani et al. (2017) - Attention Is All You Need
    \item Devlin et al. (2019) - BERT: Pre-training of Deep Bidirectional Transformers
    \item Liu et al. (2019) - RoBERTa
    \item Sanh et al. (2019) - DistilBERT
    \item Dao et al. (2022) - FlashAttention
\end{itemize}

\vspace{0.5em}

\textbf{Cognitive Neuroscience:}
\begin{itemize}\setlength{\itemsep}{0pt}
    \item Hagoort \& Indefrey (2014) - The neurobiology of language beyond single words
    \item Kuperberg \& Jaeger (2016) - What do we mean by prediction in language comprehension?
    \item Willems et al. (2016) - Prediction during natural language comprehension
\end{itemize}

\vspace{0.5em}

\textbf{Tutorials:}
\begin{itemize}\setlength{\itemsep}{0pt}
    \item HuggingFace Course: Chapters 1, 7
    \item The Illustrated Transformer: \url{https://jalammar.github.io/illustrated-transformer/}
    \item BertViz: \url{https://github.com/jessevig/bertviz}
\end{itemize}
\end{frame}

% ============================================================
% LOOKING FORWARD COURSE
% ============================================================
\begin{frame}[shrink=10]{Looking Forward in the Course 🔮}
\textbf{Where do we go from here?}

\vspace{1em}

\textbf{Upcoming Topics:}
\begin{itemize}\setlength{\itemsep}{0pt}
    \item \textbf{Week 7:} Decoder models and text generation (GPT family)
    \item \textbf{Week 8:} Scaling laws and large language models
    \item \textbf{Week 9:} Prompting, in-context learning, and instruction tuning
    \item \textbf{Week 10:} Alignment, RLHF, and ethical considerations
\end{itemize}

\vspace{1em}

\textbf{The Journey Continues:}
\begin{itemize}\setlength{\itemsep}{0pt}
    \item From understanding (BERT) to generation (GPT)
    \item From supervised learning to few-shot learning
    \item From narrow tasks to general-purpose models
    \item From academic research to societal impact
\end{itemize}

\vspace{1em}

\begin{center}
\Large
\textbf{The transformer revolution continues! 🚀}
\end{center}
\end{frame}

% ============================================================
% QUESTIONS
% ============================================================
\begin{frame}{Questions? 🙋}
\begin{center}
\LARGE
\textbf{Discussion Time}
\end{center}

\vspace{1em}

\textbf{Topics to discuss:}
\begin{itemize}\setlength{\itemsep}{0pt}
    \item BERT applications
    \item Understanding vs. pattern matching
    \item Cognitive neuroscience connections
    \item Limitations and future work
    \item Assignment 4 questions
\end{itemize}

\vspace{1em}

\begin{center}
\Large
Thank you! 🙏

\vspace{1em}

See you in Week 7 for GPT and text generation!
\end{center}
\end{frame}

\end{document}
